%! TEX program = lualatex
% For exporting as word use:
% pandoc -s sheep-arc-main.tex --filter pandoc-crossref --citeproc --csl nature.csl --bibliography=sheep-arc-refs.bib -o ../Liu_et_al_R1.docx

\documentclass[a4paper, hidelinks, 11pt]{article}

%%%%%%%%%%%%%%%%% LOAD PACKAGES %%%%%%%%%%%%%%%%
\usepackage[british]{babel}
\usepackage{microtype}

% font settings
\usepackage{fontspec}
\usepackage{mathtools}
\usepackage{textcomp}
\usepackage{eurosym}
\usepackage{unicode-math}

\defaultfontfeatures{Ligatures=TeX}

\setmainfont[
    Extension=.otf,
    UprightFont= *-regular,
    BoldFont=*-bold,
    ItalicFont=*-italic,
    BoldItalicFont=*-bolditalic,
    ]{TeX Gyre Pagella}

\setsansfont[
        Extension=.otf,
        UprightFont= *-regular,
        BoldFont=*-bold,
        ItalicFont=*-italic,
        BoldItalicFont=*-bolditalic,
        ]{TeX Gyre Heros}

\setmathfont{TeX Gyre Pagella Math}

\setmonofont{TeX Gyre Cursor}[Scale=1]


% headings
\usepackage{titlesec}

\titleformat*{\section}{\large\bfseries\sffamily\MakeUppercase}
\titleformat*{\subsection}{\normalsize\bfseries\sffamily}
\titleformat*{\subsubsection}{\normalsize\bfseries\sffamily}
\titleformat*{\paragraph}{\normalsize\bfseries\sffamily}
\titleformat*{\subparagraph}{\normalsize\bfseries\sffamily}

% page setup
\usepackage{fancyhdr}
\usepackage{authblk}
\usepackage{booktabs} % For prettier tables
\usepackage{color}
\usepackage{lineno}
% \modulolinenumbers[5]
\usepackage{csquotes}
\linenumbers

% Text alignment (comment out for fully justified)
\usepackage[document]{ragged2e}

% prevent widows and orphans
\usepackage[defaultlines=2,all]{nowidow}
\usepackage{needspace}

% create outline and links in pdf
\usepackage{hyperref}
\usepackage[numbered]{bookmark}
\makeatletter
\renewcommand\@seccntformat[1]{}
\makeatother

% debug warnings
% \overfullrule=1cm


%%%%%%%%%%%%%%%%% TABLES & FIGURES %%%%%%%%%%%%%%%%
% captions
\usepackage[labelfont=footnotesize,textfont=footnotesize,singlelinecheck=off,justification=raggedright]{subcaption}
\DeclareCaptionLabelSeparator{pipe}{\textrm{ | }}
\usepackage[figurename=Fig.,
            font={sf,footnotesize},
            labelfont=bf,
            labelsep=quad]{caption}

% figures:
\usepackage{graphicx}
\usepackage{wrapfig}
\usepackage{floatrow}

% tables:
\usepackage{tabulary}
\usepackage{multirow}
\usepackage{ltablex}
\usepackage{longtable} % Note: It may be necessary to compile the document several times to get a multi-page table to line up properly
\usepackage[table,xcdraw]{xcolor}
\usepackage{threeparttable}

%%%%%%%%%%%%%%%%% BIBLIOGRAPHY %%%%%%%%%%%%%%%%
\usepackage[
    style=nature,
    sorting=none,
    backend=biber,
    ]{biblatex}
\addbibresource{sheep-arc-refs.bib}
\renewcommand*{\bibfont}{\sffamily\footnotesize}

%%%%%%%%%%%%%%%%% ANNOTATIONS %%%%%%%%%%%%%%%%
\usepackage{soul} % highlighting & strikeout
\setlength{\marginparwidth}{2cm} % make room for comments
\usepackage[colorinlistoftodos, textsize=tiny]{todonotes}

%%%%%%%%%%%%%%%%% PAGE SETUP %%%%%%%%%%%%%%%%
\topmargin0.3cm
\oddsidemargin0.0cm
\textwidth15cm
\textheight22cm
\footskip1.3cm

\pagestyle{fancy}
\fancyhead{}
\fancyhead[RO]{\emph{Teladorsagia circumcincta} vaccinomics}
% \fancyfoot[RO, LE]{\thepage{}}
\fancyfoot[c]{\thepage{}}

\setlength{\headheight}{14pt}
\renewcommand{\headrulewidth}{0.4pt}
\renewcommand{\footrulewidth}{0pt}

%%%%%%%%%%%%%%%%% TITLE PAGE %%%%%%%%%%%%%%%%

\title{Vaccine-induced time- and age-dependent mucosal immunity to gastrointestinal parasite infection}

\author[1,\dag]{Wei Liu} % 0000-0001-9953-2813
\author[2,\dag,*]{Tom N. McNeilly} % 0000-0001-6469-0512
\author[2]{Mairi Mitchell}
\author[2]{Stewart T.G. Burgess} % 0000-0002-9268-7641
\author[2]{Alasdair J. Nisbet} % 0000-0002-4965-5396
\author[2,3,\dag]{Jacqueline B. Matthews} % 0000-0001-9405-5560
\author[1,2,\dag,*]{Simon A. Babayan} % 0000-0002-4949-1117

\affil[1]{Institute of Biodiversity, Animal Health and Comparative Medicine, University of Glasgow, Glasgow, G12 8QQ, Scotland, UK,}
\affil[2]{The Moredun Research Institute, Pentlands Science Park, EH26 0PZ, Scotland, UK,}
\affil[3]{Roslin Technologies Limited, Roslin Innovation Centre, University of Edinburgh, Easter Bush, EH25 9RG, Scotland, UK}
\affil[ \dag]{These authors contributed equally to this work}
\affil[*]{Corresponding authors: tom.mcneilly@moredun.ac.uk, simon.babayan@glasgow.ac.uk}

% \renewcommand\Authfont{\normalsize}
\renewcommand\Affilfont{\normalsize}

% Include the date command, but leave its argument blank.
\date{}

% Set title for the abstract (default is "Abstract")
\addto\captionsbritish{\renewcommand\abstractname{\vspace{-\baselineskip}}}

%%%%%%%%%%%%%%%%% END OF PREAMBLE %%%%%%%%%%%%%%%%

\begin{document}
% Import list of figures
% FIGURE LIST

%%%%%%% MAIN FIGURES %%%%%%%

\newcommand{\FigParasites}{
    \begin{wrapfigure}{R}{0.5\textwidth}
       \vspace{-16pt}
        \ffigbox[\FBwidth]
                {
                    \includegraphics[width=0.93\textwidth, keepaspectratio=true]{sheep-arc-figures/Fig 1 parasite counts.pdf}
                }
                {
                    \caption{\textbf{Vaccine- and age-mediated control of \emph{T. circumcincta} infection.} \textbf{a}, Worm burdens and \textbf{b}, total egg output in 3 month-old lambs with vaccination (3mo-Vax), 3 month-old lambs with adjuvant only (3mo-Ctrl), 6 month-old lambs with vaccination (6mo-Vax), and 6 month-old lambs with adjuvant only (6mo-Ctrl). Immunisation with the prototype vaccine led to a median 73.7\% reduction (P\textsubscript{vacc}~$=$~0.002, GLM) in worm counts and 50\% reduction (P\textsubscript{vacc}~$=$~0.026, GLM) in cumulative egg output.}
                     \label{fig:parasites}
                 }
       \vspace{-12pt}
    \end{wrapfigure}
}

% Alternative Fig 1
% \newcommand{\FigParasites}{
%     \begin{figure}[hp]

%         \includegraphics[width=0.6\textwidth, keepaspectratio=true]{sheep-arc-figures/Fig 1 parasite counts.pdf}

%         \caption{\textbf{Vaccine- and age-mediated control of \emph{T. circumcincta} infection.} \textbf{a}, Worm burdens and \textbf{b}, total egg output in 3 month-old lambs with vaccination (3mo-Vax), 3 month-old lambs with adjuvant only (3mo-Ctrl), 6 month-old lambs with vaccination (6mo-Vax), and 6 month-old lambs with adjuvant only (6mo-Ctrl). Immunisation with the prototype vaccine led to a median 73.7\% reduction (P\textsubscript{vacc}~$=$~0.002, GLM) in worm counts and 50\% reduction (P\textsubscript{vacc}~$=$~0.026, GLM) in cumulative egg output.}

%         \label{fig:parasites}

%     \end{figure}
% }

\newcommand{\FigCircle}{
\begin{figure}[hp]
    \includegraphics[width=\textwidth, keepaspectratio=true]{sheep-arc-figures/Fig 2 pathway circle.pdf}

    \caption{\textbf{Temporal dynamics of the expression of immune
        pathways that predict either post-mortem worm burden or
        cFEC identified by supervised machine learning
        ElasticNet.} Heat map indicates the pathway enrichment
        reported as negative $log_{10}^{\mathrm{P-value}}$,
        where the dark red denotes stronger enrichment within
        each pathway. Each concentric circle represents one of
        six time-points at which abomasal biopsies were taken
        after challenge infection. Worm burdens predicted by these
        pathways were measured at D49 post challenge.}

    \label{fig:circular}
\end{figure}

}

\newcommand{\FigEnetcorr}{
    \begin{figure}[htbp]
    \includegraphics[width=\textwidth, keepaspectratio=true]{sheep-arc-figures/Fig 3 Enet_corr.pdf}

    \caption{\textbf{Heat map of correlation coefficients} between
        \textbf{a}, worm burdens or \textbf{b}, cFEC and gene
        expression in the selected pathways before challenge
        (D0), seven (D7) and 21 (D21) days post challenge in
        vaccinated and control lambs. Colours represent negative
        (red) and positive (blue) Pearson correlation
        coefficients for each comparison. The immune pathways
        were clustered using k-means according to their
        correlation patterns over time in the vaccinated group.}

    \label{fig:enet_corr}
\end{figure}

}

\newcommand{\FigWGCNA}{
    \begin{figure}[htbp]
    \includegraphics[width=\textwidth, keepaspectratio=true]{sheep-arc-figures/Fig 4 WGCNA_gene_expression.pdf}

    \caption{\textbf{Gene expression within immune pathways
    significantly affected by immunisation, age, or both.}
        \textbf{a} Heat map of gene expression over time within
        the four treatments (3mo-Vax, 3mo-Ctrl, 6mo-Vax, and
        6mo-Ctrl). High to low scaled expression is denoted with
        blue to white hues. \textbf{b} Heat map of a generalised
        mixed model t value statistic (number of standard
        deviations from the mean)~\supercite{Bates2015} indicating
        the effect of age, vaccination, and the interaction
        between age and vaccination on the expression of each
        pathway. Purple indicates a dampening effect on the
        expression of the pathway while yellow indicates
        increased expression of that pathway relative to the
        global average. Pathways included are those presented in
        Fig.~\ref{fig:circular}. The full list of gene expression
        is available in Fig.~\ref{suppfig:PathwayDynamics}.}

    \label{fig:WGCNA_gene_expression}
\end{figure}
}

%%%%%%% SUPPLEMENTARY FIGURES %%%%%%%

\newcommand{\SuppFigTimeline}{
    \begin{figure}[H]
        \includegraphics[width=\textwidth, keepaspectratio=true]{sheep-arc-figures/Fig S1 exp timeline.pdf}

       \caption{\textbf{Experimental design.} Timeline of immunisation, challenge, and sampling of sheep.}

        \label{suppfig:timeline}
    \end{figure}
}

\newcommand{\SuppFigFlow}{
    \begin{figure}[H]
        \includegraphics[width=\textwidth, keepaspectratio=true]{sheep-arc-figures/Fig S2 analysis pipeline.pdf}

        \caption{\textbf{Data analysis pipeline overview.} Depiction of the
            flow through which raw transcript counts were taken
            to extract immunological information relevant to the
            response to vaccination and infection.
        }

        \label{suppfig:flow}
    \end{figure}
}

\newcommand{\SuppFigTSNE}{
    \begin{figure}[H]
    \includegraphics[width=\textwidth, keepaspectratio=true]{sheep-arc-figures/Fig S3 t_sne.pdf}

    \caption{\textbf{Unsupervised clustering of transcriptomes per
        treatment and day post challenge.} The $\sim$15K
        transcripts were reduced to two components using
        t-Distributed Stochastic Neighbour Embedding
        (t-SNE)~\supercite{Maaten2008} to visualise the datasets in
        two-dimensional space. Samples clustered weakly by
        treatment but more distinctly by day post challenge
        (DPC), specifically before \emph{vs.} after 14 DPC.}

    \label{suppfig:t-sne}
\end{figure}
}

\newcommand{\SuppFigEnetmodules}{
\begin{figure}[H]
    \includegraphics[width=\textwidth, keepaspectratio=true]{sheep-arc-figures/Fig S4 ElasticNet coefficients.pdf}

    \caption{\textbf{ElasticNet coefficients bar plot of WGCNA modules
        associated with worm burden and cFEC across 42 days post
        challenge.} Only modules with ElasticNet coefficients
        ≠ 0 are depicted and were retained for further
        analysis.
    }

    \label{suppfig:enetmodules}
\end{figure}
}

\newcommand{\SuppFigDiffexp}{
    \begin{figure}[H]
        \includegraphics[width=\textwidth, keepaspectratio=true]{sheep-arc-figures/Fig S5 diff_exp.pdf}

        \caption{\textbf{Time dynamics of differentially expressed (DE)
            genes.} \textbf{a}, Number of DE genes in each WGCNA
            modules within each of 4 pairwise DE pairwise
            comparisons: 3 month-old control \emph{vs.}
            6 month-old control (`3mo-Ctrl vs 6mo-Ctrl'),
            3 month-old vaccinated \emph{vs.} 3 month-old
              control (`3mo-Vax \emph{vs} 3mo-Ctrl'),
              6 month-old vaccinated \emph{vs.} 6 month-old
                control (`6mo-Vax \emph{vs} 6mo-Ctrl') and
              3 month-old vaccinated \emph{vs.} 6 month-old
                vaccinated (`3mo-Vax \emph{vs} 6mo-Vax').
                \textbf{b}, Venn diagram of all genes from four
                pairwise comparisons. \textbf{c} Heat map of the
                temporal expression patterns of significant DE
                gene expression. Individual genes within
                clusters 1—7 are shown in columns (see full
            gene list in Supplementary Data 2). Each row
        represents either 3mo or 6mo vaccinated lambs at
    different time points.}

        \label{suppfig:diff_exp}
    \end{figure}
}

\newcommand{\SuppFigPathwayDynamics}{
\begin{figure}[H]
    \includegraphics[width=\textwidth, keepaspectratio=true]{sheep-arc-figures/Fig S6 pathway_dynamics.pdf}

    \caption{\textbf{Time course of pathways significantly represented
        in abomasal transcriptomes.} Mean scaled expression
        levels in each pathway predictive of worm burdens and/or
    cFEC, selected in ~\nameref{fig:circular}. Points represent individual lambs at each time point and lines represent corresponding mean values for: vaccinated (solid lines); control (dashed lines); 3-month-old (red); and 6-month-old (blue) lambs.}

    \label{suppfig:PathwayDynamics}
\end{figure}
}


% Double-space the manuscript.
\baselineskip24pt

% Make the title.
\maketitle

% \listoftodos

\begin{abstract}

\noindent\bf\textsf{Individuals vary broadly in their response to vaccination and subsequent challenge infection, with poor vaccine responders causing persistence of both infection and transmission in populations. Yet despite having substantial economic and societal impact, the immune mechanisms that underlie such variability, especially in infected tissues, remain poorly understood. Here, to characterise how antihelminthic immunity at the mucosal site of infection developed in vaccinated lambs, we inserted gastric cannulae into the abomasa of three-month- and six-month-old lambs and longitudinally analysed their local immune response during subsequent challenge infection. The vaccine induced broad changes in pre-challenge abomasal immune profiles and reduced parasite burden and egg output post-challenge, regardless of age. However, age affected how vaccinated lambs responded to infection across multiple immune pathways: adaptive immune pathways were typically age-dependent. Identification of age-dependent and age-independent protective immune pathways may help refine the formulation of vaccines, and indicate specificities of pathogen-specific immunity more generally.}

\end{abstract}

\textbf{Keywords:} Helminth; vaccine; immunoparasitology; machine learning; pathway analysis; ageing.

\pagebreak
\section{Introduction}

Individuals vary widely in their responses to vaccination and little is known
about the causes underlying low vaccine efficacy, nor how effective a vaccine
will be ``in the real world'' against the pathogen for which it has been
developed. In fact, for most pathogens, imperfect vaccination is the rule rather
than the exception~\supercite{Weinberg2010}. This is particularly problematic
for helminths, as vaccination is \textcolor{red}{one of} the most promising alternative\textcolor{red}{s} to anthelmintic
drug treatments~\supercite{Bartsch2016} in the face of extensive drug
resistance, which is rising in human parasites~\supercite{Campbell2016,
Doyle2017} and ubiquitous in many helminths of livestock, including in
sheep~\supercite{Redman2015, Doyle2018,RoseVineer2020}. In both vaccinated and
non-vaccinated animals, even a minority of infected individuals can release
substantial numbers of helminth eggs into the environment, ensuring that
infection persists within populations~\supercite{Lloyd-Smith2005}. While the
need for better prediction of vaccine responsiveness has long been recognised,
the complexity of the factors involved — from variation in genetic
background~\supercite{Beraldi2007a, Periasamy2014} to differences in age, sex,
and immune history~\supercite{Pedersen2011, Babayan2011}, combined with immune
evasion strategies deployed by helminths to ensure persistent
infection~\supercite{McNeilly2013, Finlay2014, Grencis2015} — have so far
hampered the development of effective sub-unit vaccines for controlling
helminths~\supercite{Babayan2012a}. Resistance to several anthelmintic compounds
in sheep nematodes in the UK alone is estimated to cost in excess of £40M
annually~\supercite{Charlier2020}. With the need for effective vaccines ever
more pressing, it is therefore urgent to identify which immune pathways mediate
vaccine efficacy, and to understand how they change over time and with host age.

We recently developed a prototype subunit vaccine against \emph{Teladorsagia
circumcincta}, a major contributor to parasitic gastroenteritis in sheep. This
parasitic nematode resides in the abomasum (the gastric compartment of the
ruminant stomach) and is primarily a cause of disease in lambs. Third-stage
larvae (L3) penetrate glands within 24 h of infection and grow rapidly,
undergoing two moults before emerging into the abomasal lumen approximately 10
days post-infection~\supercite{McNeilly2009}. The resulting pathology manifests
as anorexia, diarrhoea and poor productivity. The prototype vaccine achieves
58–70\% reduction in worm burdens and up to 73–92\% fewer eggs at the peak of
worm egg shedding compared to challenge controls~\supercite{Nisbet2013}. Such
reductions in cumulative worm faecal egg counts (cFEC) are expected to have a
substantial impact on pasture contamination and modelling studies indicate that
reductions of this magnitude will parallel effective anthelmintic-driven
parasite control measures~\supercite{Barnes1995}. However, substantial variation
in the efficacy of this vaccine has been observed among
individuals~\supercite{Nisbet2013, Nisbet2019}. Such variability can be due to
multiple genetic and environmental factors~\supercite{Stear1997, Nussey2014,
Grencis2015}. While genetic factors can be manipulated by selective breeding and
environmental factors such as nutritional resource can be optimised under
farming conditions, any vaccine against \emph{T. circumcincta} is likely to be
most effective if administered to lambs before exposure to worms on pasture and
to protect them during the late spring/summer period in which their growth rate
is most susceptible to the impact of parasitic gastroenteritis. This, however,
raises concerns since the immature immune system is known to respond poorly to
immunisation~\supercite{McClure1998, McRae2015, Babayan2018b}. We therefore
sought to compare the efficacy and the immune responses to our prototype vaccine
in two different age groups, 3-month-old (3mo) and 6-month-old (6mo) lambs to
attempt to understand age- and vaccine-dependent protective immunity at the site
of infection to inform further optimisation of the vaccine.

How lambs of different ages respond immunologically to our prototype vaccine
over the course of repeated exposure and during chronic infection by \emph{T.
circumcincta} is not known. We therefore analysed the immune gene expression
pathways elicited by immunisation and subsequent repeated challenge infection
with \emph{T. circumcincta} in the abomasum using sequential gastric biopsies
aligned to a novel machine-based learning approach, to characterise: (a) which
immune pathways are associated with variation in parasite burdens and parasite
egg output and when they are expressed during the course of infection; (b) at
what time during the course of the infection are those pathways associated with
protection, and (c) how those pathways are affected by both vaccination and
host age.

\section{Results}

\subsection{Vaccination reduced worm burdens and egg shedding in
both 3mo and 6mo lambs.}

We sought to identify the immune mechanisms underlying variation in the ability
of the prototype vaccine to control \emph{T. circumcincta} infection in lambs
and to characterise how lamb age affects vaccine efficacy. Three- and
six-month-old lambs were vaccinated using our previously-published protocol
~\supercite{Nisbet2013}. Briefly, fifteen 3mo and fifteen 6mo lambs were
administered an eight-protein cocktail of \emph{T. circumcincta} recombinant
antigens with Quil-A adjuvant, three times at three-weekly intervals. All
vaccinated and non-vaccinated age-matched control lambs were then all exposed
to repeated (``trickle'') infections with 12 consecutive challenges with
\emph{T. circumcincta} infective larvae (L3) delivered in measured doses
spanning four weeks to mimic field challenge conditions (Fig.~\ref{suppfig:timeline}).

% \FigParasites

\textcolor{red}{At time of necropsy, control and vaccinated lambs had median worm burdens of 3,425 and 900, respectively, and median cumulative faecal egg counts of 7680 and 3839, respectively.} Vaccination \textcolor{red}{therefore} led to significant reductions in both worm burden assessed at
\emph{post mortem} (\nameref{fig:parasites}A; \textcolor{red}{73.7\% reduction,} P~\(=\)~0.002) and cumulative
faecal egg counts (cFEC) (\nameref{fig:parasites}B; \textcolor{red}{50\% reduction,} P~\(=\)~0.026) in both age
groups relative to control lambs. In addition, 6mo lambs, whether immunised or
not, controlled worm burdens more effectively than 3mo lambs (P~\(=\)~0.012).
However, total egg output differed only marginally between age groups and was
characteristically variable.

\subsection{Pathway enrichment was most predictive of vaccine efficacy prior
to challenge and following parasitological events.}

Because worm burdens and cFEC varied substantially between individuals
regardless of age or vaccination, we then sought to identify (i) which immune
pathways were associated with controlling worm burdens at post-mortem (49 days
after the start of challenge infection) and cFEC \textcolor{red}{across all groups with vaccination and age treated as explanatory variables}, and (ii) when their
expression in the abomasum best predicted these measures of parasite infection.
To address both questions, we built an analysis pipeline to identify the immune
pathways that best predicted infection outcomes (Fig.~\ref{suppfig:flow}).
RNAseq was performed on biopsies taken repeatedly from the abomasum at temporal
intervals. These produced six transcriptomes spanning 42 days for each lamb,
beginning post-vaccination but shortly before the first challenge and ending
one week prior to the end of the trial when the worm burden analysis was
undertaken (Fig.~\ref{suppfig:timeline}). T-distributed stochastic neighbour
embedding (t-SNE) suggested there was weak structure within the gene expression
read counts of all 180 whole transcriptomes (Fig.~\ref{suppfig:t-sne}). We then
clustered transcripts into co-expression modules regardless of age using a
weighted correlation network analysis (WGCNA~\supercite{Langfelder2008}). This
generated a reduced list of variables, each of which represented a set of genes
that behaved similarly across all transcriptomes. To identify which of the
resulting modules contained genes of interest, we used worm burdens and cFEC as
response variables in ElasticNet regression models to which we mapped the
eigenvalue of each WGCNA module, and used the coefficients learned by the
ElasticNet to rank the modules according to the strength and direction of their
association with either worm burden or cFEC (Fig.~\ref{suppfig:enetmodules}).

We then extracted all genes included within the modules that predicted either
worm burden or cFEC, then entered that gene list into a canonical pathway
analysis, and selected the resulting immune pathways~\supercite{Kramer2014}.
~\nameref{fig:circular} presents the immune pathways that robustly predicted
worm burdens of cFEC and their level of enrichment (proportion of genes within
specific pathways that was detected in the transcriptome) at each time-point.
This revealed that most immune pathways that predicted either parasite burden
49 days (D49) post-infection or cFEC were already enriched by D0, i.e.
following immunisation, but prior to the first challenge. Pathways that were
significantly represented exclusively at this time-point included B cell
development and activation signalling, as well as innate pathways involved in
interferon and TNF signalling, macrophage stimulating protein receptor
signalling (MSP-RON), and the complement system (\nameref{fig:circular}, inner
ring).

% \FigCircle

All other pathways significantly enriched at D0 were also enriched at
subsequent time-points. In particular, pathways whose expression correlated
with worm or egg counts showed significant enrichment at days 7, 21, and 42
post-challenge, consistent with the timing of three significant phases of the
life cycle, i.e. at D7, upon emergence of the first cohort of late fourth stage
larvae (L4)/early adults from the gastric glands into the
lumen~\supercite{Denham1969, McNeilly2017}, coinciding with high enrichment scores
of pathways involved in the regulation of innate responses (NF-κB, eosinophil
CCR3, GMCSF, NO and ROS expression by macrophages, IL-2/IL-15, IL-3, and IL-8
signalling) as well as antigen-processing and lymphoid cell activation (FLT3,
IL-7, CD40, CXCR4); at D21, co-incident with the initiation of egg-laying by
the founding population of \emph{T. circumcincta}, a further set of immune
pathways was activated, including inflammatory and stress response signals
(LPS-stimulated MAPK, MIF, glucocorticoid receptors, IL-17, and TLR signalling)
and suppression of lymphocyte proliferation (anti-proliferative role of TOB in
T cell signalling). Enrichment at D42 was elevated in the majority of pathways
depicted in ~\nameref{fig:circular} (see full timecourse for all pathways in
~Fig.~\ref{suppfig:PathwayDynamics} and~\nameref{supdata:ipa}). Being close to
the post-mortem time-point (D49), this likely reflects near-contemporaneous and
potentially shorter-lived correlations between immune gene expression and
parasite burdens, and were thus not considered further in the prediction of
protective immunity. The very low enrichment scores in the remaining
time-points likely indicates either low expression of the corresponding genes
due to biological regulatory processes, or too much variance for any
statistical pattern to be detected. We therefore decided to further focus on
time-points D0, D7, and D21 for their relevance to predicting vaccine-mediated
immunity against \emph{T. circumcincta} in these lambs.

\subsection{Direction and strength of correlation between immune pathway
expression and parasitological measures varied between time-points.}

To identify the time-points at which the selected pathways
(\nameref{fig:circular}) best predicted reduced parasite burdens and cFEC, we
assessed the direction and strength of the correlations between each pathway
and either parasite burden (\nameref{fig:enet_corr}A) or cFEC
(\nameref{fig:enet_corr}B) before (D0) or 7 and 21 days following the first
exposure to \emph{T. circumcincta} L3. For ease of interpretation, we then
clustered immune pathways according to whether they were negatively associated
with parasite burdens at each, or all, of these time-points.

Comparing the expression of immune pathways at D0 in vaccinated lambs and
non-vaccinated controls indicated that IL-6, Th1, PPAR, and interferon
signalling pathways elicited by the vaccine were instrumental in reducing
parasite numbers at post-mortem (D49), while higher activation of IL-17A, IL-9,
CCR3, and TGF-β signalling pathways in vaccinated lambs prior to infection
predicted lower egg shedding. In non-vaccinated lambs, no pathways were
predictive of lower worm numbers, whereas IL-17A, CCR3, and chemokine
signalling pathways were identified as predictors of low cFEC.

At D7, Th1 and PPAR signalling remained negatively associated with post-mortem
worm burden in vaccinates, while most other pathways were either positively or
not associated with worm burdens. With regard to egg shedding, in vaccinated
lambs activation of IL-17A remained negatively associated with cFEC, with
pathways associated with regulation of T cell activation and Th1 polarisation
(TOB and IL-12 pathways, respectively) also associated with reduced cFEC.
Similar to D0, no immune pathways were predictive of lower parasite numbers in
control lambs, although CCR3 and chemokine signalling pathways remained
consistently, though weakly, negatively associated with cFEC.

Finally, at D21 vaccinated lambs differed from controls by exhibiting strong
negative associations between worm burdens and chemokine signalling, IL-6,
NFκB, natural killer cells, glucocorticoid receptor and acute phase response
signalling, suggesting a potential role for responses to tissue injury in
vaccine-induced protection against \emph{T. circumcincta}
(\nameref{fig:enet_corr}). Additionally, pathways associated with T cell
activation and polarisation (IL-12, DC maturation, TOB, NFAT, Th1 and Th2
signalling pathways) were all negatively associated with worm burden at this
time-point. In the non-vaccinated lambs, no pathways at D21 showed any
association with worm burdens. With regard to predicting vaccine-induced
impacts on cFEC, many of the pathways associated with reduced worm burdens at
D21, in particular chemokine signalling, IL-6, and acute phase response
signalling, were also predictive of lower cFEC in vaccinates. Lower egg
shedding in vaccinated lambs was also associated with increased expression of
CD27 in lymphocytes, and macrophage migration inhibitory factor (MIF).
Interestingly, in non-vaccinated lambs, IL-7, IL-9, FLT3, and CCR3 signalling
were negatively associated with parasite egg shedding at D21, whereas these
signalling pathways were positively associated with cFEC in vaccinated lambs at
the same time-point post-challenge.

% \FigEnetcorr

\subsection{Protective immune pathways were differentially affected by age
and vaccination status.}

We then sought to analyse how the age of lambs at immunisation affected their
expression of protective immune pathways at the same time-points. Overall, the
mean gene expression levels for most of the pathways described above were most
strongly expressed in the older (6mo) vaccinates at D21 than at the two other
time-points (\nameref{fig:WGCNA_gene_expression}A). Six-month-old lambs also
displayed a greater differential in gene expression between D0 and D21,
regardless of vaccination status. Indeed, while mean gene expression of immune
pathways over all time-points was often greater in 3mo lambs, gene expression
levels varied little over time in the 3mo lambs, particularity in
non-vaccinated individuals. To assess how age and immunisation explained the
observed variation in the expression of the focal immune pathways, we
constructed generalised linear models for each pathway in turn with age,
vaccine status, and the interaction between age and vaccine as explanatory
variables. Most protective pathways identified by the ElasticNet
(\nameref{fig:enet_corr}) were affected by lamb age, especially pathways
involved in the activation of adaptive immunity and its maintenance, spanning
Th1, Th2, and Th17 pathways (\nameref{fig:WGCNA_gene_expression}B `Age ×
Vaccine' column). Pathways that were elicited by vaccination independently of
age included acute phase response signalling, LPS-mediated inhibition of RXR
function, MIF Regulation of Innate Immunity, and Toll-like Receptor Signalling
(\nameref{fig:WGCNA_gene_expression}B, `Vaccine' column). Finally, regardless
of vaccination, only Antiproliferative Role of TOB in T Cell Signalling and
possibly the Complement System were affected by age, with 6mo lambs expressing
them at lower levels than 3mo lambs (\nameref{fig:WGCNA_gene_expression}B
`Age' column).

% \FigWGCNA

\section{Discussion}

We have recently developed a prophylactic vaccine that shows great promise
against \emph{T. circumcincta}~\supercite{Nisbet2013}. However, for this
prototype to deliver consistent protection a better understanding of the immune
mechanisms that underlie the observed imperfect immunisation in sheep of
relevant ages is needed ~\supercite{Barnes1995}. We monitored immune responses
of vaccinated sheep at the site of infection (the abomasum) over the first 42
days of a trickle challenge infection, and measured parasite burdens on day 49.
To assess how host age affects vaccine efficacy, we immunised lambs at 3 or 6
months of age. Consistent with our previous reports~\supercite{Nisbet2013,
Nisbet2016}, the vaccine led to significant reductions of both worm burdens and
egg outputs compared to adjuvant-only recipients (\nameref{fig:parasites}),
\textcolor{red}{confirming the specific protective effect of the antigens used
in the vaccine.} However, 3-month-old lambs remained more heavily infected than
6-month-old lambs throughout the experiment, consistent with previous reports in
ruminants that the immature immune system responds poorly to prior
immunisation~\supercite{Smith1985,Corpa2000}. Further, age differentially
affected the ability of sheep to control worm numbers and the numbers of
parasite eggs produced in their faeces: while we observed a significant
reduction of parasite burden in 6-month-old animals irrespective of their
vaccination status, age alone did not significantly affect total egg output
(\nameref{fig:parasites}). This suggests that maturation of the immune system
allows better control of \emph{T. circumcincta} worm burden with age but may
have a more limited effect on parasite transmission via parasite egg shedding.
Such apparent compensation and density-dependent fecundity in parasites have
been previously observed in this~\supercite{Bishop2000} and other host-parasite
systems~\supercite{Churcher2005,Churcher2006,Bleay2007}. Further, while the
prototype vaccine reduced parasite and egg densities, it did not totally
eliminate the presence of high worm egg shedders, raising the possibility that
inherently ``wormy'' individuals may not respond well to vaccination. Reducing
infection in these poor vaccine responders is of particular importance for
limiting the spread of infection in vaccinated flocks. \textcolor{red}{In
addition, while Quil A has been reported to affect \emph{Fasciola
hepatica}~\supercite{Hacariz2009}, future studies will need to test whether and
how Quil A alone may affect GI nematodes including \emph{T. circumcincta}.}

In this study, we took a novel systems vaccinology approach to identify immune
transcriptional pathways that predict both post-mortem worm burdens and cFEC in
vaccinated individuals. While similar approaches have been taken to understand
the mechanisms by which vaccines stimulate protection, these have largely
focused on repeated analysis of blood leukocyte transcriptomes which are
clustered into blood transcriptional modules (BTM) and subsequently correlated
with antibody or cellular immune phenotypes~\supercite{Li2014,Hou2017,
Braun2018}. This approach, while providing important information on immune
pathways involved in vaccine responses, relies on capturing transcriptomic
signatures of recirculating leukocytes which may not truly reflect the immune
response at the site(s) of infection or vaccination. In contrast, the
application of a gastric cannulation technique in this study allowed repeated
sampling of the gastric mucosa, resulting in an unparalleled temporal evaluation
of the local transcriptomic responses over time at the parasite's predilection
site.  Furthermore, we used a novel machine-learning approach to maximise the
generalisability of the qualitative and quantitative associations between immune
responses and worm numbers or fecundity. Thus, the combination of repeated
\emph{in situ} sampling, unbiased selection of immune mediators of
vaccine-driven protection, and pathway analysis, allowed us to generate a
biologically-interpretable, robust, and dynamical representation of how
vaccination affects the immune response to \emph{T. circumcincta} infection in
lambs.

Using this approach, analysis of the immune response at the site of infection
revealed that three main events were indicative of interactions between
vaccination and parasite life history and were predictive of parasite burdens.
First, priming of the immune system by the vaccine prior to challenge: immune
pathways that determined the worm burdens measured at day 49 post challenge were
already highly activated in the transcriptomes before challenge, likely driven
by the pre-challenge vaccination. Second, the response to parasite life history
events: while reduced pathway enrichment observed between days 7 and 28 suggests
a down-regulation of immune activity in the abomasum once founder populations of
parasites had established infection, we observed an activation of more directed
responses to ongoing parasitological events involved in innate responses (e.g.
eosinophil CCR3, NF-κB activation and NO/ROS expression by macrophages) and
antigen-processing and lymphocyte recruitment and activation (CD40, IL-7, CXCR4
signalling) 1–2 weeks post-challenge, matching the timing of the emergence of
the first larvae from the gastric glands~\supercite{Simpson2000}.  This was
followed by inflammatory and stress responses (e.g. IL-17 signalling,
TLR-signalling, glucocorticoid receptor signalling) around 3 weeks after
infection when parasites begin producing eggs (see ~\nameref{fig:circular}).
And third, the tighter correlations between gene expression levels and parasite
numbers at the last sampling time-point immediately prior to post-mortem
analysis of worm burdens (day 42) point to confounding between the effects of
vaccination and the temporal proximity between immune tissue sampling (day 42)
and measurement of parasite counts (day 49).

Further investigation therefore focused on three time-points, days 0, 7 and 21,
which were most predictive of parasite numbers and/or fecundity, to identify
pathways that were associated with vaccine-induced protection. Of most interest
were protective pathways enriched in the mucosa of vaccinated lambs immediately
prior to challenge. These included Th1, IL-6 and interferon signalling pathways
which negatively correlated with worm burdens, and IL-17A, IL-9, CCR3 and TGF-β
pathways associated with lower parasite egg output, and suggested that the
systemically delivered vaccine was able to modulate the immune system at a
distant mucosal site. This was a rather surprising observation, as while immune
signatures have previously been reported in peripheral blood, for example
following yellow fever and influenza
vaccination~\supercite{Kotliarov2020,Hou2017}, systemically delivered vaccines
are generally considered poor at inducing mucosal immune responses due to the
inability of the vaccines to induce appropriate homing receptors on activated
lymphocytes~\supercite{Clements2016}. Furthermore, while some adjuvants (e.g.
TLR agonists and bacterial ADP-ribosylating toxin adjuvants) have been shown to
confer mucosal homing properties~\supercite{Frederick2018}, this has not been
widely reported for saponin-based adjuvants such as Quil-A used here, although
our results are consistent with an earlier study of a systemically delivered
\emph{Ostertagia ostertagi} subunit vaccine formulated with Quil-A that reported
similar mucosal priming of immune cells, and in particular natural killer (NK)
cells~\supercite{Gonzalez-Hernandez2016}.

Regardless of the mechanism by which this mucosal priming operates, the vaccine
appears to promote a protective Th1/Th17 type response within the mucosa, with
evidence of active Th17 polarisation potentially via IL-6 and
TGF-β~\supercite{Li2007d,Mangan2006,Zhou2007a}, as well as evidence of enrichment of
Th2-associated pathways, CCR3 and IL-9, more established effectors of
anti-parasite immunity~\supercite{Faulkner1997}, possibly via MIF, which was
recently reported as essential to type 2 immunity against \emph{Heligmosomoides
polygyrus} in mice~\supercite{Filbey2019}. Interestingly, some of these pathways
(MIF, IL17A, and CCR3) were also associated with reduced parasite egg output in
control lambs, suggesting that of the protective pathways induced by the
vaccine, anti-worm Th1 responses were most unique. The association of
vaccine-induced protection with Th1 immunity was maintained at day 7
post-infection, with Th1 pathways being associated with both worm burdens and
egg output. By 21 days, a wider range of protective pathways were identified,
including those associated with tissue injury and inflammation\textcolor{red}{, which might have been driven by ongoing tissue damage caused by worms as consecutive cohorts progressed through their life cycle}. At this
time-point, both Th1 and Th2 signalling pathways were associated with
protection, indicating a broadening of the immune response.

While these associations will require further causal validation, they suggest
that for optimal protection, the vaccine may prime the mucosa towards a Th1 and
potentially Th17-type response early in infection before broadening out to a
more mixed Th1/Th2 response. This is potentially contentious, as it is
well-established in numerous animal models that protective immunity to
gastrointestinal parasites, including \emph{T. circumcincta}, is associated
with Th2 immunity, whereas Th1 and Th17 immunity is associated with
susceptibility~\supercite{Craig2007,Venturina2013,Grencis1997,Maizels2009a,Gossner2012a}.
However, this appears to be a rather simplistic model based on
counter-regulation of Th immune responses derived from tightly-controlled
laboratory models and/or analysis of responses at limited time-points
post-infection. Indeed, our longitudinal sampling at the site of parasite
infection coupled with whole-transcriptome analysis is likely to have allowed a
finer description of the sequence of protective responses, revealing
contributions to the protective response of Th1, Th17 and Th2 responses at
different phases of the infection within the same individual. These results are
consistent with a previous study in sheep, in which natural resistance to
\emph{T. circumcincta} was associated with early Th1 responses prior to
development of Th2 immunity~\supercite{Hassan2011a}, and a more recent study in
which early activation of the Th17 pathway was associated with natural
resistance to \emph{Haemonchus contortus} in goats~\supercite{Aboshady2020}. The
role of Th17 in protection is potentially explained by the ability of this
pathway to elicit innate lymphoid cells and multipotent progenitor type 2 cells
early in infection that subsequently promote CD\(4^+\) Th2 cells and associated
cytokine expression~\supercite{Saenz2010,Saenz2013,Smith2018}. This then activates
antiparasitic effector cells such as
eosinophils~\supercite{Martin2000,Hernandez2020}, consistent with our finding that
CCR3 signalling in eosinophils at D21 post challenge was negatively associated
with egg output (\nameref{fig:enet_corr}). The previous association of Th17
responses with susceptibility was determined after more long-standing (12 week)
\emph{T. circumcincta} infection~\supercite{Gossner2012a}, where the association
could be explained by triggering of Th17 responses secondary to gut barrier
disruption~\supercite{Kempski2017}.

Further investigation into how vaccination and age affected the expression of
these immune pathways revealed that 6mo lambs increased the expression of the
pathways to a greater level than did 3mo lambs at D21, which coincides with the
first emergence of parasites from gastric crypts into the lumen of the
abomasum~\supercite{Denham1969,McNeilly2017}. This age-dependant expression of
protective immune pathways largely involved the activation and maintenance of
the adaptive response. This late maturation of the adaptive response to
helminths observed in lambs is consistent with reports of later maturation of T
cells in sheep~\supercite{Outteridge1985,Colditz1996} and in other host-parasite
systems as previously reported~\supercite{Babayan2018b}. Conversely, innate pathways
appeared age-independent for these two cohorts. However, vaccination alone
significantly explained the upregulation of innate pathways such as those
involved in acute phase response signalling, LPS/IL-1 mediated inhibition of
RXR functions, TLR signalling, and MIF regulation pathways, only a subset of
which are reported to be enhanced by saponin-based
adjuvants~\supercite{Ho2018,Marty-Roix2016}, the others likely induced by the
infection.

In conclusion, most protective-associated pathways were induced by vaccination
and in older animals; the immune pathways found to control adult worms and eggs
only partially overlapped and were activated at different phases of the
challenge infection, indicating the need for anthelmintic vaccines to stimulate
a broad set of pathways, rather than just antibody production alone; protective
pathways enriched pre-challenge in vaccinates suggest specific adjuvants, such
as those promoting Th1/Th17 responses may be useful to improve vaccine
performance.


\begingroup
% Use smaller sans font
\sffamily
\small
% Double-space the methods section.
\baselineskip24pt

\section{Methods }

\subsection{Lambs \& Infection}

Texel cross lambs \textcolor{red}{reared under conditions to exclude helminth
infection (confirmed by FEC analysis)} were randomly allocated to four groups
which were balanced for weight and sex. \textcolor{red}{Two lambs were removed
from the experiment early due to health concerns.} The \textcolor{red}{final}
groups comprised vaccinated 3-month-old lambs (n~\(=\)~15 \textcolor{red}{[7
females, 8 males]}), control (adjuvant only) 3-month-old lambs (n~\(=\)~16
\textcolor{red}{[8 females, 8 males]}), vaccinated 6-month-old lambs (n~\(=\)~15
\textcolor{red}{[8 females, 7 males]}), and control (adjuvant only) 6-month-old
lambs (n~\(=\)~16 \textcolor{red}{[8 females, 8 males]}). All lambs were
infected with 2,000 \emph{T. circumcincta} L3 stage larvae three times per week
for four weeks beginning on the day of the final immunisation.

\subsection{Prototype vaccine formulation}

Each lamb in the vaccinated groups was injected with 400 µg of a recombinant
protein mix as described previously~\supercite{Nisbet2013}, containing 50µg of
each protein. PBS soluble proteins Tci-ASP-1, Tci-MIF-1, Tci-TGH-2,
Tci-APY-1, Tci-SAA-1, Tci-CF-1 and Tci-ES20 were administered as a single
injection with 5mg Quil-A (Brenntag Biosector). Tci-MEP-1 was administered
separately in PBS with 2M urea and 5mg Quil-A. Injections were given
subcutaneously at two sites on the neck. \textcolor{red}{To test the antigen-specific effects of the vaccine, c}ontrol animals received injections
containing PBS/urea and Quil-A only. Lambs received three immunisations at
three-weekly intervals (day 0, 21, 42) with the first immunisation at 3 or 6
months of age (Fig.~\ref{suppfig:timeline}).

\subsection{Sampling}

Seven animals in each vaccine group and eight animals in the control groups
were fitted with abomasal cannulae, as previously described~\supercite{Huntley2004},
to allow repeated biopsy of the abomasal mucosa throughout the trickle
infection period. Briefly, peri-operative analgesia was administered prior to
anaesthetic induction using Meloxicam (Metacam®, Boehringer Ingleheim) at 1
mg/kg body weight (BW). Anaesthesia was induced by intra-venous Propofol
(PropoFlo™, Zoetis) at 3mg/kg BW and maintained using Isoflurane (IsoFlo™,
Zoetis). Following surgical preparation of the site, the abomasum was located
and exteriorised. Abomasal cannulae, constructed from a modified disposable
10 ml syringe barrels (PlastiPak™, Becton and Dickinson)~\supercite{Huntley2004},
were fitted midway between the mesenteric border and the greater curvature of
the lateral wall, approximately 7 cm cranial to the pylorus. The free end of
the cannulae were exteriorised through a laparotomy incision in the abdominal
wall, then anchored using external neoprene flanges. Surgery to fit the
cannulae was performed between the second and third immunisation time-points.
Over 49 days, abomasal biopsies were taken using a pair of 30 cm long 5 mm ×
2 mm punch mucosal biopsy forceps (Richard Wolf GmBH), inserted via the
cannulae as follows: three biopsies per animal at each time-point on days 0,
7, 14, 21, 28, 42, and 49 (Fig.~\ref{suppfig:timeline}). The three biopsies
per animal were taken in a clockwise manner to ensure different sites of the
abomasal mucosa were sampled. Samples were placed immediately into RNAlater
(Sigma-Aldrich) and stored at -80 °C for subsequent RNA extraction.

\subsection{RNA-Seq library preparation and sequencing}

Total RNA was extracted from the abomasal biopsies as follows: the 645
abomasal biopsy samples were homogenised in RLT buffer (Qiagen Ltd, UK) using
a Precellys bead basher (Bertin Instruments, UK) with CK28 bead tubes
(Stretton Scientific, UK). Samples were centrifuged at 14,000~g at 4°C for 10
mins and the supernatant collected for processing using a RNeasy
mini-isolation kit (Qiagen Ltd, UK) according to the manufacturers’ protocol,
including an on-column DNase digestion. RNA quality and integrity were
assessed using a Nanodrop spectrophotometer (Thermo Fisher, UK) and a
Bioanalyser RNA Nanochip (Agilent Technologies Ltd, UK). The yield of total
RNA was determined on a Qubit Fluorometer (Thermo Fisher, UK) using the Broad
Range RNA kit (Thermo Fisher, UK). The RNA isolated from three biopsies per
animal for each time-point was pooled 1:1:1 by weight to generate the final
samples for RNA-seq assessment. \textcolor{red}{All samples exceeded RIN ≥ 8.2.} The resulting 180 RNA samples were sequenced
on an Illumina NextSeq 500 by Glasgow Polyomics~\supercite{PolyomicsURL}
generating 75 bp paired-end reads at an average sequencing depth of
25~Mbp/sample. \textcolor{red}{FastQC analysis indicated that all sequences exceeded Phred Scores of 30 along all 75bp.}

\subsection{RNA-Seq quality control and alignment}

The sequencing quality controls were finished by FastQC (v0.11.5), which
provides a comprehensive report for each RNA-Seq sample. Base calls were made
using the Illumina CASAVA 1.8 pipeline. All 180 samples passed the QC filters
(MultiQC~\supercite{Ewels2016} report, \nameref{supdata:multiqc}), suggesting that
RNA extraction and subsequent sequencing were of good quality and cutadapt
(v1.11) was used for adapter trimming. Pseudo alignment of the read data to
the latest version of the sheep transcriptome (cDNA) (Oar-v3.1) was performed
with Kallisto v0.46.2~\supercite{Bray2016}, generating read count data for each
transcript across all samples. The R package tximport~\supercite{Soneson2015} was
used to prepare the abundance matrix for downstream analysis. All samples
were normalised using DESeq2~\supercite{Love2014} and genes with read counts
greater than five across at least two samples were selected for downstream
analysis.

\subsection{Gene network analysis and machine learning}

We first evaluated whether any structure in gene expression could be
visualised using t-Distributed Stochastic Neighbour Embedding (t-SNE). To
reduce the dimensionality of the transcriptome, we then used the R package
Weighted Correlation Network Analysis (WGCNA)~\supercite{Langfelder2008} to
generate gene co-expression networks from vaccinated groups (3mo-Vax and
6mo-Vax) at six time-points after start of the challenge (Days 0 – 42), using
the full transcriptome. WGCNA groups genes and builds networks using the
co-expression similarity measure defined as \(
S_{i,j}=\left(\frac{1}{2}+\frac{1}{2}corr\left(x_i, x_j\right)\right)^\alpha
\), where \(S_{i,j}\) is the correlation between gene expressions \(x_i\) and
\(x_j\), and \(\alpha\) is the soft threshold weight selected by scale-free
topology criterion~\supercite{Zhang2005,Langfelder2008} set at 6 (Day 0), 8 (Day
7), 7 (Day 14), 12 (Day 21), 8 (Day 28), and 9 (Day 42). The eigengene of
each cluster was used to quantify its overall expression. To select the genes
which best predicted the parasitological read-out of interest (i.e., worm
burdens and cFEC), we used the ElasticNet algorithm~\supercite{Zou2005a} from
Python’s Scikit-Learn software library to fit a linear regression between the
eigengenes and worm burdens or cFEC, and then ranked the gene clusters by
their resulting coefficients. All WGCNA modules for which the ElasticNet
coefficient was not null were retained for further pathway analysis.

\subsection{Pathway analysis}

Pathway enrichment was generated with Ingenuity Pathway Analysis (IPA, QIAGEN
Inc.)~\supercite{Kramer2014} in which each gene identifier was mapped to its
corresponding gene object in Ingenuity’s Knowledge Base using canonical
pathway analysis to identify the biological pathways of most significance.
Only immune pathways among those identified in IPA (\nameref{supdata:ipa})
were retained for further analysis.

\subsection{Time series differential gene expression analysis}

The R package maSigPro~\supercite{Nueda2014}, a two-step regression to find
significant differences between treatments over time, was used for gene
differential expression analysis with multiple time-points (\nameref{supdata:diffexp}).

\subsection{Statistical analysis of age and vaccination effects on parasite counts and gene pathway expression}

The effects of age and immunisation on both worm burdens and cFEC were assessed
with generalised linear models for negative binomial distributions with a log
link, and residuals tested for normality using the Jarque-Bera test for
normality. The effects of age, immunisation, and their interaction on the
eigengene of each immune pathway were assessed with linear mixed models using
sampling date adn sheep ID as a random effects to account for repeated sampling.
The t statistic was used to indicate effect sizes as fold deviation between
group means~\supercite{Bates2015}.

\subsection{Ethics statement}

All animal procedures were performed at Moredun Research Institute (MRI) under
Home Office licence 70/7914. Ethical approval was obtained from the MRI Animal
Welfare and Ethical Review Body (E22/15).

\section{Data availability}

All links to transcriptomes and code for bioinformatics, machine learning, and
statistical analysis are available at \url{https://github.com/SimonAB/Liu2021}.

\section{Acknowledgements}

We thank Leigh Andrews, Alison Morrison and Dave Bartley, Moredun Research
Institute, for their help in conducting animal studies and in the provision of
parasite material and the Bioservices Unit, Moredun Research Institute, for
expert care of the animals. We also thank Manus Graham and Roy Davie, Moredun
Research Institute, for assistance with gastric cannulation surgeries.
\textcolor{red}{The authors thank two anonymous reviewers for their insightful
and constructive feedback.} This work was supported by the Biotechnology and
Biological Sciences Research Council (Grant references BB/M011968/1 and
BB/M012956/1 to JBM and SAB, respectively).

\section{Competing Interests}
The authors declare that there are no competing interests.

\section{Author contributions}

TMN, STGB, AJN, JBM, and SAB conceived of and designed the study; TNM and MM performed the experiments; WL and SAB performed the analyses; TMN, STGB, AJN, JBM, and SAB wrote the manuscript. \textcolor{red}{WL and TNM contributed equally as co-first authors; JBM and SAB contributed equally as senior authors.}


\endgroup

\section{References}
\printbibliography[heading=none]

\section{Figure Legends}

\subsection{Fig. 1} \label{fig:parasites}

\textbf{Vaccine- and age-mediated control of \emph{T. circumcincta} infection.} \textbf{a}, Worm burdens and \textbf{b}, total egg output in 3 month-old lambs with vaccination (3mo-Vax), 3 month-old lambs with adjuvant only (3mo-Ctrl), 6 month-old lambs with vaccination (6mo-Vax), and 6 month-old lambs with adjuvant only (6mo-Ctrl). Immunisation with the prototype vaccine led to a median 73.7\% reduction (P\textsubscript{vacc}~$=$~0.002, GLM) in worm counts and 50\% reduction (P\textsubscript{vacc}~$=$~0.026, GLM) in cumulative egg output.



\subsection{Fig. 2} \label{fig:circular}

\textbf{Temporal dynamics of the expression of immune
        pathways that predict either post-mortem worm burden or
        cFEC identified by supervised machine learning
        ElasticNet.} Heat map indicates the pathway enrichment
        reported as negative $log_{10}^{\mathrm{P-value}}$,
        where the dark red denotes stronger enrichment within
        each pathway. Each concentric circle represents one of
        six time-points at which abomasal biopsies were taken
        after challenge infection. Worm burdens predicted by these
        pathways were measured at D49 post challenge.

\subsection{Fig. 3} \label{fig:enet_corr}

\textbf{Heat map of correlation coefficients} between
        \textbf{a}, worm burdens or \textbf{b}, cFEC and gene
        expression in the selected pathways before challenge
        (D0), seven (D7) and 21 (D21) days post challenge in
        vaccinated and control lambs. Colours represent negative
        (red) and positive (blue) Pearson correlation
        coefficients for each comparison. The immune pathways
        were clustered using k-means according to their
        correlation patterns over time in the vaccinated group.

\subsection{Fig. 4} \label{fig:WGCNA_gene_expression}
\textbf{Gene expression within immune pathways
    significantly affected by immunisation, age, or both.}
        \textbf{a} Heat map of gene expression over time within
        the four treatments (3mo-Vax, 3mo-Ctrl, 6mo-Vax, and
        6mo-Ctrl). High to low scaled expression is denoted with
        blue to white hues. \textbf{b} Heat map of a generalised
        mixed model t value statistic (number of standard
        deviations from the mean)~\supercite{Bates2015} indicating
        the effect of age, vaccination, and the interaction
        between age and vaccination on the expression of each
        pathway. Purple indicates a dampening effect on the
        expression of the pathway while yellow indicates
        increased expression of that pathway relative to the
        global average. Pathways included are those presented in
        \nameref{fig:circular}. The full list of gene expression
        is available in Fig.~\ref{suppfig:PathwayDynamics}.

\pagebreak

\newcommand{\sectionbreak}{\clearpage}

\section{Supplementary information}
\setcounter{page}{1}
\renewcommand{\thepage}{S\arabic{page}}
\setcounter{table}{0}
\renewcommand{\thetable}{S\arabic{table}}
\setcounter{figure}{0}
\renewcommand{\thefigure}{S\arabic{figure}}

\subsection{Supplementary Data 1}\label{supdata:multiqc}
\texttt{multiqc-report-biopsy.html}: Quality check report of 180 RNA-seq sample FastQC reports using MultiQC~\supercite{Ewels2016}.

\subsection{Supplementary Data 2}\label{supdata:ipa}
\texttt{ipa\_3V\_6V.ods}: Ingenuity Pathway analysis (IPA) results for all the WGCNA clusters shown in ~\nameref{fig:circular}.

\subsection{Supplementary Data 3}\label{supdata:diffexp}
\texttt{gene-cluster-masigpro.ods}: Differentially expressed gene list between the age groups. The expression levels are shown in ~Fig.~\ref{suppfig:diff_exp}C.

\subsection{Supplementary Figures}
\newcommand{\subsectionbreak}{\clearpage}

\subsection{Study design}

\SuppFigTimeline

\subsection{Methods flow overview}

\SuppFigFlow

\subsection{Clustering of all transcriptomes using t-SNE.}

\SuppFigTSNE

\subsection{Weight of WGCNA modules predicting parasite burdens.}

\SuppFigEnetmodules

\subsection{Differential expression analysis of genes included in WGCNA modules.}

\SuppFigDiffexp

\subsection{Time-course of immune pathways significantly
represented in abomasal transcriptomes}

\SuppFigPathwayDynamics

\end{document}
